
El grupo de investigación \href{http://arantxa.ii.uam.es/~gnb/}{\textit{GNB}}, cuenta con una serie de plugins para \textit{RTXI}, que extraen algunas de las funcionalidades de \textit{RTHybrid} \cite{Reyes-Sanchez2020, Amaducci2019} y hacen uso de ellas en \textit{RTXI}. Dichos plugins presentan un inconveniente a la hora de usarlos en versiones superiores a la 3.0.0, pues en esta actualización la clase en la que se basan, `DefaultGuiModel`, ha sido retirada por completo \cite{RTXI2023}. Esto hace que no se puedan compilar ni usar con las nuevas versiones de \textit{RTXI}.

Por ello, se han migrado todos los plugins de la anterior versión a la nueva (Repositorio \ref{SEC:RTHPRTXI3REPO}). Para realizar dicha migración, se ha usado el script `create\_template\_plugin.sh`, que genera una plantilla para el plugin. Con esta base se han ido completando los plugins reciclando el código de la anterior versión, pero bajo las nuevas clases de `widgets.cpp`, que dividen la funcionalidad del módulo en `Panel` (para la GUI), `Component` (para la lógica) y `Plugin` (para el propio \textit{RTXI}).

La adversidad más remarcable durante la implementación ha sido mostrar las variables en la GUI, pues los `STATES` que se usaban anteriormente, ahora solo pueden ser enteros sin signo, por lo que se han tenido que usar colas y un temporizador para ir pasando los valores de interés a la GUI. Esta idea, se ha obtenido analizando los componentes base de \textit{RTXI}, en concreto el \textit{RT Benchmarks} (localizado en el directorio `plugins/performance\_measurement/` de \textit{RTXI}).
